% Generated draft: requires mathematical and source review.
\documentclass[11pt]{article}
\usepackage[margin=25mm]{geometry}
\usepackage{fontspec}
\setmainfont{Latin Modern Roman}
\usepackage{amsmath,amssymb,mathtools,unicode-math}
\setmathfont{Latin Modern Math}
\usepackage{enumitem,needspace,xurl,hyperref}
\setlength{\parindent}{0pt}
\setlength{\parskip}{5pt}
\setlength{\emergencystretch}{3em}
\newcommand{\N}{\mathbb{N}}
\newcommand{\Z}{\mathbb{Z}}
\newcommand{\Q}{\mathbb{Q}}
\newcommand{\R}{\mathbb{R}}
\newcommand{\C}{\mathbb{C}}
\newcommand{\F}{\mathbb{F}}
\newcommand{\A}{\mathbb{A}}
\newcommand{\PP}{\mathbb{P}}
\newcommand{\E}{\mathbb{E}}
\newcommand{\eps}{\varepsilon}
\newcommand{\dd}{\,\mathrm{d}}
\begin{document}
\section*{UnsolvedMath ID: 30006522}
\textbf{Source problem code:} catalogue entry
\subsection{Definitions and mathematical statement}
Let \(n\geq 1\), and let \(S_n\) be the symmetric group on \(\{1,\ldots,n\}\). For \(1\leq k<n\), let \(s_k=(k\;k+1)\). For \(v\in S_n\), define its Coxeter length by
\[
\ell(v)=\#\{(i,j):1\leq i<j\leq n,\ v(i)>v(j)\}.
\]
The left weak order \(\leq_L\) is the transitive closure of the relations \(v<_L s_kv\) with \(\ell(s_kv)=\ell(v)+1\). Thus left multiplication here interchanges the values \(k\) and \(k+1\).

Under the conventional interpretation using ordinary single Grothendieck polynomials, let \(G_w\) denote the Grothendieck polynomial associated with \(w\), and let \(\mathfrak S_v\) denote the Schubert polynomial associated with \(v\). Write its Schubert expansion as
\[
G_w=\sum_{v\in S_n}a_{w,v}\mathfrak S_v,
\qquad
\operatorname{Supp}_S(G_w)=\{v\in S_n:a_{w,v}\neq0\}.
\]
Here \(\deg(G_w)\) means total polynomial degree.

\textbf{Conjecture 2, as supplied, under this finite-group interpretation.} For every integer \(n\geq1\), every \(w\in S_n\), and every \(u\in\operatorname{Supp}_S(G_w)\),
\[
\ell(u)<\deg(G_w)
\quad\Longrightarrow\quad
\exists k\in\{1,\ldots,n-1\}\quad
\bigl(u<_L s_ku\ \text{ and }\ s_ku\in\operatorname{Supp}_S(G_w)\bigr).
\]
Equivalently, the required step satisfies \(\ell(s_ku)=\ell(u)+1\). When \(n=1\), the displayed index set is empty and the quantified implication retains its usual logical meaning.

\textbf{Unresolved conventions.} The supplied excerpt does not explicitly define \(G_w\), its normalization, or \(\operatorname{Supp}_S\), nor does it specify the ambient group for support indices or the range of \(k\). The definitions above make the usual finite, single-polynomial interpretation explicit; they require confirmation against the complete report. If support is instead indexed in \(S_\infty=\bigcup_{m\geq1}S_m\), using the standard embeddings that fix additional letters, the Schubert expansion and support must be indexed there and the unrestricted simple-reflection index is \(k\geq1\). The excerpt alone does not explicitly resolve this convention.

The surrounding source text additionally reports that Conjectures 1 and 2 together would imply analogous conjectures labelled \([3]\), with Schubert support replaced by monomial support and left weak order replaced by monomial divisibility. It mentions a stronger statement implying Conjecture 1 and a subsequent conjecture about left weak-order intervals. Neither those statements nor the identification of \([3]\) is supplied, so these contextual claims are retained as source reports rather than additional formal assertions.
\subsection{Short English statement}
The supplied conjecture asks whether every permutation in the Schubert support of \(G_w\) whose length is below \(\deg(G_w)\) admits a length-increasing left multiplication by an adjacent transposition that remains in the support. The finite-group formulation uses conventional polynomial and support definitions that require confirmation from the full report. The supplied status is partially solved; its dated assessment says that structured families are known while the arbitrary-permutation case remains unresolved.
\subsection{Sources}
\begin{itemize}
\item Local UnsolvedMath registry: supplied complete source record for permanent ID \(30006522\), titled \emph{Weak-Order Growth of Grothendieck Schubert Supports}, exact source problem code \texttt{catalogue entry}, in the Oberwolfach Reports Open Problems collection. The record gives proposed year 2026 and no named proposer.
\item Supplied report citation: \emph{Enumerative Combinatorics} (2026), Oberwolfach Reports, DOI \texttt{10.4171/owr/2026/2}, pp.~81--156, \url{https://doi.org/10.4171/owr/2026/2}. The excerpt labels the target Conjecture 2; no exact page or individual contribution attribution is supplied.
\item Supplied literature source: Mészáros--Setiabrata--St. Dizier, \emph{On the support of Grothendieck polynomials}, \url{https://arxiv.org/abs/2201.09452}.
\item Supplied literature source: Chou--Yu, \emph{Grothendieck polynomials of inverse fireworks permutations}, \url{https://arxiv.org/abs/2403.12186}.
\item Literature assessment date, exactly as supplied: \texttt{2026-08-22}. The supplied dated curation labels the problem partially solved and states that support conjectures and stronger variants are proved for several structured families, while the arbitrary-permutation statement remains unresolved. This is supplied dated curation derived from the OWR record, not an independently verified literature assessment; time-sensitive status and cited claims require independent verification.
\item The registry labels the statement \texttt{corrected\_verified}, explaining that only formula delimiters, notation, grammar, and typography were repaired and that mathematical content and scope were preserved. This supplied editorial label is not independent mathematical verification. The excerpt's reference \([3]\) is not identified in the supplied record.
\end{itemize}
\end{document}
